\documentclass[12pt,a4paper]{report}

\usepackage[utf8]{inputenc}
\usepackage[T1]{fontenc}
\usepackage{lmodern}
\usepackage{geometry}
\usepackage{setspace}
\usepackage{graphicx}
\usepackage{booktabs}
\usepackage{tabularx}
\usepackage{array}
\usepackage{enumitem}
\usepackage{xcolor}
\usepackage{hyperref}
\usepackage{titlesec}
\usepackage{fancyhdr}
\usepackage{natbib}
\usepackage{url}

\geometry{
  left=30mm,
  right=30mm,
  top=28mm,
  bottom=30mm
}

\onehalfspacing
\setlength{\parindent}{0pt}
\setlength{\parskip}{0.65em}
\setlist[itemize]{topsep=0.3em,itemsep=0.15em,leftmargin=1.4em}
\setlist[enumerate]{topsep=0.3em,itemsep=0.25em,leftmargin=1.6em}

\definecolor{ChalmersBlue}{HTML}{004B93}
\definecolor{SoftBlue}{HTML}{EAF2FA}
\definecolor{TextGray}{HTML}{222222}

\hypersetup{
  colorlinks=true,
  linkcolor=ChalmersBlue,
  citecolor=ChalmersBlue,
  urlcolor=ChalmersBlue
}

\titleformat{\chapter}
  {\normalfont\huge\bfseries\color{ChalmersBlue}}
  {\thechapter.}{0.6em}{}
\titlespacing*{\chapter}{0pt}{-10pt}{18pt}

\titleformat{\section}
  {\normalfont\Large\bfseries\color{ChalmersBlue}}
  {\thesection}{0.6em}{}
\titlespacing*{\section}{0pt}{14pt}{6pt}

\titleformat{\subsection}
  {\normalfont\large\bfseries\color{ChalmersBlue}}
  {\thesubsection}{0.6em}{}
\titlespacing*{\subsection}{0pt}{10pt}{4pt}

\pagestyle{fancy}
\fancyhf{}
\lhead{\small Role-Play Driven UX}
\rhead{\small RECS Toolkit Report}
\cfoot{\thepage}
\renewcommand{\headrulewidth}{0.4pt}

\newcolumntype{Y}{>{\raggedright\arraybackslash}X}
\renewcommand{\arraystretch}{1.25}

\begin{document}

% ------------------------------------------------------------------ Title page
\begin{titlepage}
  \centering
  \vspace*{35mm}
  {\Huge\bfseries\color{ChalmersBlue} Role-Play Driven UX\par}
  \vspace{8mm}
  {\Large A TTRPG-Inspired Toolkit for User Needs Analysis\par}
  \vspace{18mm}
  {\large Boyang Wei \quad Yi Yang\par}
  \vspace{4mm}
  {\normalsize Interaction Design, Chalmers University of Technology\par}
  \vspace{6mm}
  {\normalsize Master thesis open project\par}
  \vfill
  {\large \today\par}
\end{titlepage}

% ------------------------------------------------------------------ Front matter
\pagenumbering{roman}

\chapter*{Abstract}
\addcontentsline{toc}{chapter}{Abstract}

User needs analysis depends heavily on what participants are willing and able to say.
When research concerns emotionally complex or cognitively demanding products, direct
self-report often falls short.
This project explores whether structures from tabletop role-playing games (TTRPGs) can
extend the reach of UX needs research.
Building on the Requirements Elicitation Campaign Setting (RECS) proposed by
\citet{sturdee2023ttrpg}, we designed a working research toolkit that translates TTRPG
concepts---characters, scenarios, dice mechanics, and GM facilitation---into a practical
instrument for early-stage UX work.
The central design contribution is a conceptual reframing of the dice mechanic: rolls do
not judge product success; they inject perturbations into the research environment,
keeping the real product and the participant's actual behaviour at the centre of
observation.
Poco, the first author's master thesis prototype for ADHD-oriented productivity support,
served as the primary test vehicle.
The result is a session-complete toolkit that supports character creation, scenario
selection, perturbation injection, GM note-taking, quote capture, and structured export.

\clearpage
\tableofcontents
\clearpage

\pagenumbering{arabic}

% ================================================================== Chapter 1
\chapter{Introduction}

UX research often depends on asking people to describe what they do, what they need,
and what frustrates them.
But needs are not always accessible through direct reflection.
Some only become visible when a person is placed inside a situation, asked to make
decisions under pressure, and allowed to respond to tension, interruption, or
uncertainty.
This is especially relevant for products that support emotional, cognitive, or everyday
self-management---where the user's experience is shaped not only by interface structure,
but also by attention, energy, emotional state, and context.

This problem is particularly acute for products designed for users with ADHD or
attention difficulties.
Such users may find it hard to articulate what slows them down without being in the
moment of slowdown itself.
Abstract task completion metrics or retrospective interview questions often miss the
texture of these experiences.

This project asks whether ideas from tabletop role-playing games can create that situated
moment within a research session.
In a TTRPG, participants do not simply answer questions---they take on roles, act within
a fictional situation, and respond to events introduced by rules and by a facilitator.
\citet{fine1983shared} frames role-playing games as social worlds built through shared
imagination, rules, and group interaction.
\citet{bowman2010functions} further argues that role-playing supports identity
exploration, problem-solving, and community formation.
Both suggest that the role-playing frame creates a kind of productive distance:
participants can speak and act from within a character, which may make it easier to
express needs, hesitations, and reactions that feel too personal to state directly.

The direct inspiration for this project is \citet{sturdee2023ttrpg}, who propose that
UX research can borrow TTRPG structures to create ``playable UX worlds.''
Their proposal is conceptual; it does not provide a working instrument.
This project responds to that gap.
The contribution is a practical toolkit that turns those ideas into a runnable research
session, demonstrated through Poco---a mobile-first productivity app prototype designed
to reduce the friction of starting tasks for users who experience attentional and
emotional barriers.

% ================================================================== Chapter 2
\chapter{Background}

\section{TTRPGs as structured imagination}

Tabletop role-playing games combine rules, collaborative storytelling, and facilitated
uncertainty.
A GM presents the world and its situations; players control characters within that world.
The rules provide structure, but the course of events emerges through conversation,
interpretation, and decision-making.

This balance between structure and openness is relevant to UX research.
Traditional research sessions can become over-scripted: the participant performs a task,
the researcher observes whether it succeeds, and the session moves on.
The TTRPG model suggests a different approach---the participant inhabits a role,
encounters situations through that role, and responds to events that neither participant
nor researcher can fully predict.
\citet{fine1983shared} describes this as a shared social world with its own norms and
interaction patterns.
\citet{bowman2010functions} extends this by identifying how role-playing helps
participants safely explore responses, emotions, and decisions they might otherwise avoid
naming.

For UX research, the useful parts of TTRPG are not the fantasy content.
They are the facilitation structure, character framing, scenario-based action, and
rule-governed uncertainty.

\section{RECS: TTRPG structures applied to UX}

\citet{sturdee2023ttrpg} provide the central theoretical anchor for this project.
Their paper proposes a Requirements Elicitation Campaign Setting (RECS), in which UX
research sessions are reimagined as playable worlds.
Characters relate to personas; adventures relate to user journeys; the GM role relates
to research facilitation; dice mechanics introduce rule-governed disruption.

The most important idea borrowed from RECS is psychological distance.
When a participant speaks as a character, they may be able to describe frustration,
shame, confusion, or resistance in a less personally exposed way.
For products like Poco---where the user's internal state is part of the design
problem---this distance can surface data that direct questioning cannot.

The RECS paper is deliberately speculative.
It proposes a direction but does not specify how to run a session, what rules to use,
or what data to collect.
This project treats that openness as a design brief.

\section{Distinction from gamification}

This project is often misread as gamification, so the distinction is worth stating
directly.
\citet{deterding2011gamification} define gamification as the use of game design
elements in non-game contexts to influence motivation and behaviour.
Gamification is typically about incentives---points, badges, leaderboards, progress bars.

This toolkit is not designed to motivate participants or reward them.
The dice mechanic is not entertainment.
The character sheet is not a levelling system.
The GM role is not dramatic performance.
These elements exist because they create the structural conditions under which situated
behaviour becomes observable and recordable.
The toolkit's game-like surface is in service of a research function, not a motivational
one.

% ================================================================== Chapter 3
\chapter{Design Process}

\section{Design goals}

The toolkit was designed around six requirements, derived from the needs of a practical
research instrument and from the gaps in the RECS proposal.

\begin{enumerate}
  \item \textbf{Support persona-as-character creation.}
        The participant should be able to build or adjust a character with attributes,
        traits, and a short backstory, making the session inhabitable rather than only
        task-based.
  \item \textbf{Use scenario-based research tasks.}
        Sessions should unfold through a sequence of beats---alternating between
        contextual narration and product-directed challenges.
  \item \textbf{Separate product performance from dice outcomes.}
        The dice should not decide whether the participant succeeds in the product.
        Product behaviour on the real interface remains observable.
        The dice only determine whether the environment introduces a disruption.
  \item \textbf{Give the researcher a clear GM role.}
        The researcher must be able to guide narration, call for rolls, observe
        participant behaviour, capture notes, and record quotes without losing track of
        the session.
  \item \textbf{Produce analysable session records.}
        The session should conclude with a structured export containing the character
        sheet, rolls, perturbations, GM notes, quotes, and debrief.
  \item \textbf{Remain reusable across products.}
        Poco is the main test vehicle, but the toolkit should not be locked to it.
        Product-specific scenarios live in separate packs so that other researchers can
        adapt the method.
\end{enumerate}

\section{Character sheet and attribute system}

The character sheet is the participant's entry point into role-play.
Rather than testing as themselves, participants create a plausible user character with a
name, role, age, backstory, device context, and six attributes:

\begin{itemize}
  \item \textbf{Patience (PAT)} --- tolerance for friction and repetition
  \item \textbf{Focus (FOC)} --- ability to sustain directed attention
  \item \textbf{Working Memory (WMEM)} --- capacity to hold multiple things in mind
  \item \textbf{Emotional Regulation (EREG)} --- management of frustration, shame, and
        discomfort
  \item \textbf{Tech Literacy (TLIT)} --- fluency with digital interfaces
  \item \textbf{Sensory Tolerance (SENS)} --- sensitivity to visual and auditory overload
\end{itemize}

Each attribute is rated from 1 to 5, with a total point budget of 18.
A budget of 18 enables a neutral character (all attributes at 3) but encourages
deliberate trade-offs: a high-Focus character may have low Emotional Regulation; high
Tech Literacy may come with low Patience.
These trade-offs make the persona more specific and give the GM a clearer model of how
the character would respond to friction.

The attributes were chosen because they describe the user capacities that most strongly
shape interaction with products like Poco.
Task initiation, attentional sustain, cognitive overload, and emotional self-management
are not peripheral to the design problem---they are central to it.

\section{Scenario and beat structure}

Each session is built around a scenario.
A scenario has a premise, a goal, and a sequence of beats.
Narration beats set emotional and situational context through text the GM reads aloud.
Challenge beats identify a concrete product action, assign an attribute and difficulty,
and trigger a roll.

Difficulty classes are calibrated at three levels: Easy (DC~10), Medium (DC~14), and
Hard (DC~18), reflecting the complexity of the underlying product interaction.

The Poco pack contains three scenarios:

\begin{itemize}
  \item \textbf{The Morning Standoff}---morning task initiation, state check-in, and
        choosing one task to begin; probing how the app supports first contact with the
        task list.
  \item \textbf{Focus Under Siege}---setting up a focus session, responding to the Mochi
        companion, and sustaining attention under simulated distraction; probing how the
        app behaves when focus is already fragile.
  \item \textbf{The Overwhelmed Mind Drop}---capturing scattered thoughts into the inbox,
        converting one into a task, and beginning; probing how well the app accommodates
        an overloaded state.
\end{itemize}

These scenarios were selected because they map to Poco's core thesis: productivity
support should reduce the friction of starting, not only optimise output.

\section{The central design pivot: from success judgement to perturbation injection}
\label{sec:pivot}

The most consequential design decision was a reframing of the dice mechanic---one that
emerged only after identifying a conceptual error in the first version of the toolkit.

In the initial interpretation, the dice roll was treated as a success or failure
judgement: a high roll meant the character succeeded at the product interaction; a low
roll meant failure.
This reading was intuitive because nearly all TTRPG systems use dice this way.
But applied to UX research, it produced a serious problem.

If the die decides whether the participant succeeds at a product task, then the
die---not the product---becomes the evaluative agent.
The participant may respond to the roll result rather than to the actual interface.
A participant who fails a roll might avoid exploring the interface because the fiction
tells them they failed.
A participant who succeeds might move on without expressing genuine difficulty.
In either case, the researcher is no longer observing the participant's real interaction
with the real product.
The toolkit has inserted itself between the participant and the product, undermining the
research goal.

This is a departure from how the RECS framework describes the dice mechanic.
\citet{sturdee2023ttrpg} frame rolls as part of the rules-governed session structure but
do not specify what dice results should govern.
The paper maps dice broadly to ``rule-based uncertainty'' without defining whether that
uncertainty applies to participant actions or to the research environment.
This ambiguity, interpreted naively, leads to the v1 error described above.

The corrected design treats rolls as \emph{environmental perturbation signals}, not
success judgements.
A roll does not determine what the participant did with the product.
It determines what the surrounding world injects into the situation.
The participant still uses the real product, and their actual behaviour---navigation
decisions, verbal hesitations, moments of confusion, emotional reactions---remains the
research data.
The roll only modulates the environment:

\begin{table}[h]
\centering
\small
\begin{tabularx}{\textwidth}{>{\bfseries}p{0.32\textwidth}Y}
\toprule
Roll band & Environmental signal \\
\midrule
Critical success & Exceptionally calm; no perturbation, small advantage \\
Success          & Calm continuation; no perturbation \\
Mixed            & Mild perturbation injected \\
Failure          & Significant perturbation injected \\
Critical failure & Compound perturbation (multiple or severe) \\
\bottomrule
\end{tabularx}
\caption{Dice outcome bands and their environmental meaning.}
\end{table}

Under this logic, a low roll on a Focus check during a Poco session does not mean the
participant failed to use the app.
It means the environment now introduces a distraction---a notification, an intrusive
thought, a companion event, or a moment of self-doubt.
The participant continues using the real product, and the GM observes how they respond to
the combined pressure of the product and the disruption.

This reframing keeps the product at the centre of observation while allowing the
researcher to introduce realistic, controlled friction.
It also changes the relationship between dice outcomes and scenario design: every beat
outcome now has a plausible research function, rather than acting as a pass/fail gate
that short-circuits product interaction.

This shift was guided by the design iteration principle described by
\citet{fullerton2018game}: game-like systems require playtesting and revision to surface
assumptions that look correct in theory but fail in use.

\section{GM role design}

The toolkit is primarily researcher-facing.
The participant uses the real product on their own device; the GM uses the toolkit to
guide the session, introduce perturbations, and capture data.
This split keeps the participant's attention on the product rather than on the research
instrument.

The GM-facing layer includes narration prompts for each beat, task prompts specifying
what the participant should attempt in the real product, dice roll controls, perturbation
outcome descriptions, a live GM notebook, quote capture with timestamps, a roll and
perturbation log, and a summary and export screen.
This makes the toolkit not only a game-like framing layer, but also a research capture
tool---the session record is built incrementally during the session rather than
reconstructed afterwards.

% ================================================================== Chapter 4
\chapter{Design Outcome}

\section{Session flow}

The final toolkit follows a five-stage session flow:

\begin{table}[h]
\centering
\small
\begin{tabularx}{\textwidth}{>{\bfseries}p{0.22\textwidth}Y Y}
\toprule
Stage & Purpose & Main output \\
\midrule
Setup           & Participant code entry and session preparation & Session identity \\
Character       & Persona creation and attribute allocation & Character sheet \\
Scenario Picker & Product pack and scenario selection & Active scenario \\
Scenario Run    & Beat-by-beat narration, rolls, and perturbation injection & Roll log, GM notes, quotes \\
Summary         & Session review and structured export & Markdown and JSON record \\
\bottomrule
\end{tabularx}
\caption{Five-stage session flow of the toolkit.}
\end{table}

The flow is designed so a researcher can run one participant through one scenario in a
single sitting of 60--90 minutes.
Each stage gates the next on data readiness---a character sheet must be complete before
a scenario can be selected; a scenario must be active before the run screen
opens---preventing incomplete records.

\section{Dice mechanic and outcome bands}

Rolls use a d20 plus an attribute modifier derived from the character's relevant
attribute value.
The roll total is compared to the scenario's difficulty class to determine the
perturbation band described in Section~\ref{sec:pivot}.

Because the modifier range is constrained by the 18-point budget, the system avoids
characters who never encounter friction.
Even a high-attribute character will occasionally draw perturbations, reflecting the
reality that no user is entirely resistant to the conditions these attributes describe.

\section{Complication library}

When a roll calls for a perturbation, the GM draws from a complication library organised
into four categories:

\begin{itemize}
  \item \textbf{Interface}---unexpected system behaviour, unclear controls, layout
        surprises.
        \textit{Example: ``A loading banner briefly obscures the button you were about
        to tap.''}
  \item \textbf{Cognitive}---memory lapses, decision paralysis, losing track of intent.
        \textit{Example: ``You scroll back to the top of the list, unsure if you already
        chose something.''}
  \item \textbf{Emotional}---self-blame, shame, frustration, or feeling observed.
        \textit{Example: ``The priority tags start to feel like accusations. Your
        shoulders tighten.''}
  \item \textbf{External}---off-device interruptions such as notifications, sounds, or
        social pressure.
        \textit{Example: ``A message notification lights up the lock screen edge as you
        reach for the check-in.''}
\end{itemize}

Each complication has a severity rating.
Mild perturbations draw low-severity items; compound perturbations draw multiple or
high-severity items.
The Poco pack adds product-specific complications tied to the Mochi companion, status-pill
tags, and the mood/energy/focus check-in.

The four-category structure reflects the actual distribution of friction in digital
product use: users are affected by the interaction of interface behaviour, cognitive
state, emotional condition, and environmental context simultaneously.
The library gives the researcher a systematic way to introduce any of these dimensions
without improvisation.

\section{Poco as the test vehicle}

Poco is a mobile-first productivity app prototype designed as the first author's master
thesis design work.
Its core interactions include a mood/energy/focus state check-in, a task list with
breakdown support, a Mind Drop inbox for offloading scattered thoughts, focus sessions
with a companion character named Mochi, and a break-down flow for decomposing large
tasks.

Poco was a strong test vehicle for the RECS toolkit for three reasons.
First, it deals with subjective states---energy, focus, overwhelm, and emotional
resistance to starting---that are difficult to probe through task completion metrics
alone.
Second, its value depends on how it behaves under messy everyday conditions, not only
under clean laboratory use.
Third, it contains several scenario-rich moments: approaching the task list in the
morning, entering focus mode, reacting to Mochi's feedback, or dumping thoughts into an
inbox.

Using Poco also allowed the toolkit to be stress-tested against a real design artefact
with known design intentions, rather than a hypothetical product with no accessible
design rationale.

\section{Research output structure}

The session export includes: participant code; selected pack and scenario; complete
character sheet and attribute profile; each beat encountered; all roll outcomes;
complications injected; stress and UXP changes; GM live notes; participant quotes with
timestamps; and post-session debrief.

The Markdown export is structured for human reading and qualitative analysis.
The JSON export preserves the full session object for possible aggregation or
cross-session comparison.

% ================================================================== Chapter 5
\chapter{Testing and Evaluation}

\textit{[This chapter is to be completed following pilot and main study sessions.]}

% ================================================================== Chapter 6
\chapter{Discussion}

The toolkit's main value is that it makes user needs more situated.
Rather than asking a participant to describe needs in the abstract, it places a character
inside a moment: waking up late, trying to start work, losing focus, receiving a
notification, or feeling judged by a priority tag.
These moments can surface needs that are difficult to access through direct questioning
or standard task observation.

For Poco specifically, this situatedness matters because the app's design challenge is
not only whether controls are visible or flows are understandable.
It is whether the product can help when the user is tired, scattered, ashamed,
indecisive, or overloaded.
The TTRPG-inspired toolkit gives the researcher a way to explore these states without
forcing the participant to speak from personal vulnerability.
The character creates distance; the scenario keeps the session concrete.

The perturbation mechanic adds controlled messiness.
Traditional usability sessions often occur in quiet, artificial conditions.
The toolkit lets the researcher introduce realistic disruptions---a notification, a
cognitive stutter, a moment of emotional resistance---consistently and at calibrated
severity levels.
This does not simulate every real-world condition, but it provides a structured way to
observe behaviour under stress.

There are real limitations.
The method depends on facilitation skill: a GM who over-performs the fiction may distract
from the product; one who under-explains the rules may leave the participant confused
about whether they are performing, testing, or both.
Some participants will find role-play unfamiliar or uncomfortable.
Sessions take longer than standard usability tests because character creation and scenario
framing require time that direct task observation does not.

Ethically, the method requires care.
Perturbations can introduce emotional pressure, and the researcher must avoid making the
participant feel trapped or judged.
The character frame should protect the participant by providing narrative distance, not
be used to extract responses participants would otherwise withhold.
Participants should be able to pause, step out of character, or end the session at any
point, and this must be made explicit before and during the opening of the session.

This method is therefore most appropriate when the research question concerns situated
behaviour, emotional response, complex context, or early requirements exploration.
It is less appropriate when the goal is only to identify simple interface bugs, measure
narrow task completion rates, or compare performance across many participants under
standardised conditions.

% ================================================================== Chapter 7
\chapter{Conclusion}

This project designed and implemented a TTRPG-inspired UX research toolkit for user
needs analysis.
Building on the RECS framework proposed by \citet{sturdee2023ttrpg} and on related work
on role-playing games, gamification, and iterative design, the project translated TTRPG
structures into a practical research instrument that can be run in a single session.

The most important contribution is the reframing of the dice mechanic as perturbation
injection rather than success judgement.
This change is not cosmetic.
It repositions the real product as the primary object of observation and the
participant's actual behaviour as the data---rather than allowing the dice to displace or
overshadow the product interaction.
It also makes the method more honest to the RECS vision, which asks how TTRPG structures
can make UX research richer, not how they can replace direct product observation.

Poco served as the primary test vehicle, and its fit was strong: the app's focus on task
initiation, emotional friction, and attentional management created exactly the kind of
design problem this toolkit is equipped to explore.
At the same time, the toolkit was built to remain product-agnostic through a pack-based
architecture, so future researchers can apply the method to other products without
rebuilding the core session structure.

The next step is to run a comparative evaluation between RECS-style sessions and
traditional usability testing on Poco, using a within-subjects design planned for the
main thesis study.
The expected contribution is not a claim that TTRPG methods should replace existing UX
methods, but a clearer account of when and why role-play-driven research can surface
needs, tensions, and contextual behaviours that standard methods are structurally unable
to reach.

% ================================================================== References
\bibliographystyle{apalike}
\bibliography{references}

\end{document}
